% Insert into an existing LaTeX document.
% Required packages: amsmath, amssymb, amsthm.
% Define a proposition environment in the parent document if needed.
% These are standard identities and conditional bounds instantiated for the
% candidate method. They are not claims of original theory or trained results.

\section{Conditional properties of task-weighted response shaping}
\label{sec:task-response-theory}

The physical state is $x=(p,v,a)\in\mathbb R^6$. The response model uses
$z=(\hat p,\hat v,m)$, initialized by
$E_\phi(x)=(p,v,e_\phi(x))$. The latent memory $m$ is not the measured
actuator state $a$. With $\Delta t=0.05$ and $K=4$, its recursion is
\begin{align}
 \hat p^+&=\hat p+\Delta t\hat v,\qquad
 \hat v^+=\hat v+\Delta t m,\nonumber\\
 m^+&=m+\frac{\Delta t}{\tau}\bigl(h_\phi(z,c,\kappa)-m\bigr).
 \label{eq:tr-model}
\end{align}
The response $b_j\in\mathbb R^5$ contains normalized physical position,
velocity, and the realized resource $\rho_j=\|u_j\|^2/U^2$. It contains
neither $a_j$ nor $m_j$. Let $B=(b_0,\ldots,b_{K-1})$, let $Q$ be a
frozen open-loop model prediction from the same physical initial state,
and let $\mathcal W=\operatorname{diag}(W_0,\ldots,W_{K-1})$.

\paragraph{Implemented alternating soft objectives.}
For $N$ complete episodes with $M$ macro intervals each, define
\begin{align}
 L_{\mathrm{local}}&=\frac{\Delta t}{MN}\sum_{m,n,j}
       \|b_{mjn}-q_{\phi,mjn}\|_{W_{mjn}}^2,\nonumber\\
 L_{\mathrm{long}}&=\frac{\Delta t}{MN}\sum_{m,n,j}
       \|b_{mjn}-\bar q_{\phi,mjn}\|_{W_{mjn}}^2,\nonumber\\
 L_{\mathrm{model}}&=L_{\mathrm{local}}+0.25L_{\mathrm{long}}
       +\lambda_{\mathrm{task}}J_g(\phi;\bar C_{\mathrm{recorded}},\kappa_0).
 \label{eq:tr-model-objective}
\end{align}
The local prediction starts from each actual macro initial state. The
long prediction is encoded only once at the episode initial state and
uses that episode's frozen executed command sequence. The task term
uses the same complete recorded sequence and includes integrated stage
and terminal costs; it does not call MPC to generate a new plan.
The proposal uses $\lambda_{\mathrm{task}}=0.1$ only in enabled arms.
Metrics instead use each macro's online plan, extended to the remaining
task endpoint by holding its final command. The two task surrogates
share the model and task-cost definition, but need not share a command
sequence or gradient.

The actor reward, with global microstep index $t=mK+j$ and environment
index omitted, is
\begin{equation}
 r_{m,j}=-\Delta t\,\ell_{\mathrm{task}}(x_{t+1},u_t,\kappa_{t+1})
     -\Delta t\,\beta\|b_{mj}-q_{\phi,mj}\|_{W_{mj}}^2
     -\mathbf 1_{\{t=MK-1\}}
       \ell_{\mathrm{terminal}}(x_{t+1},\kappa_{t+1}),\qquad\beta=0.2.
 \label{eq:tr-shaped-reward}
\end{equation}
Batch A uses the deterministic-mean actor and supports a model proposal;
fresh batch B uses stochastic exploration under the frozen proposal model
and supplies shaped returns for PPO; fresh batch C uses the new
deterministic-mean actor and supports identification-only calibration with
$\lambda_{\mathrm{task}}=0$. Every arm uses these same collection modes.
The calibrated model is used by the next MPC iteration and targets the
deterministic deployment response. Matching in stochastic batch B still
penalizes exploration-response variance; its mean is not identified with
the deterministic response. All three batches and calibration work must
be counted.
Targets and metrics are saved at each macro start and remain fixed
through the actor batch and PPO epochs. These are soft consistency
penalties, not hard realization constraints. The alternating procedure
has not been proved to be exact coordinate descent of a single fixed
population objective. No parameter updates are run in this implementation
verification stage.

\begin{proposition}[Trace-normalized bounded task metric]
\label{prop:tr-metric}
For $v\in\mathbb R^d$ and $\alpha\ge0$, define
\begin{equation}
 A=I+\frac{\alpha}{1+\|v\|^2}vv^T,\qquad
 W=\frac{dA}{\operatorname{tr}(A)}.
 \label{eq:tr-metric}
\end{equation}
Write $r=\|v\|^2$ and $s=\alpha r/(1+r)$. Then
\begin{equation}
 \operatorname{tr}(W)=d,\qquad
 \frac{d}{d+\alpha}I\preceq W
 \preceq\frac{d(1+\alpha)}{d+\alpha}I.
 \label{eq:tr-metric-bounds}
\end{equation}
Moreover,
\begin{equation}
 v^TW^{-1}v=\frac{r(d+s)}{d(1+s)},\qquad
 W^{-1}=\frac{d+s}{d}\left(I-
 \frac{\alpha vv^T}{1+(1+\alpha)r}\right).
 \label{eq:tr-dualnorm}
\end{equation}
\end{proposition}
\begin{proof}
The matrix $A$ has eigenvalue one orthogonal to $v$, eigenvalue $1+s$
parallel to $v$, and trace $d+s$. Hence $W$ has eigenvalues $d/(d+s)$
and $d(1+s)/(d+s)$ in the respective directions. The first decreases
with $s$ and the second is nondecreasing for integer $d\ge1$.
Using $0\le s\le\alpha$ proves the bounds. Dividing $r$ by the parallel
eigenvalue proves the quadratic-form identity, with value zero at $v=0$.
The inverse formula follows by inverting the same rank-one update.
\end{proof}
For $d=5$ and $\alpha=4$, the spectral interval is $[5/9,25/9]$.
The trace is five, matching the uniform metric's average weight.
This removes a simultaneous increase in total metric weight but does not
make covariance-weighted penalties identical across different metrics.

\begin{proposition}[Mean matching and response variance]
\label{prop:tr-variance}
Fix a macro initial condition and command context $\xi$. Let $Q$ and
$\mathcal W$ be deterministic and frozen conditional on $\xi$. If $B$
has finite second moments under the policy, then
\begin{equation}
 \mathbb E_\theta[\|B-Q\|_{\mathcal W}^2\mid\xi]
 =\|\mu_\theta-Q\|_{\mathcal W}^2
  +\operatorname{tr}(\mathcal W\Sigma_\theta),
 \label{eq:tr-variance}
\end{equation}
where $\mu_\theta=\mathbb E_\theta[B\mid\xi]$ and
$\Sigma_\theta=\operatorname{Cov}_\theta(B\mid\xi)$.
\end{proposition}
\begin{proof}
Expand $B-Q=(B-\mu_\theta)+(\mu_\theta-Q)$. The cross term has zero
conditional expectation, while the centered quadratic term is
$\operatorname{tr}(\mathcal W\Sigma_\theta)$.
\end{proof}

Under common support, differentiation--integration interchange,
parameter-independent environment transitions, and a fixed latent-action
map, the conditional score identity is
\begin{equation}
 \nabla_\theta\mathbb E_\theta[A(B)\mid\xi]
 =\mathbb E_\theta\left[
 (A(B)-b_{\mathrm{base}}(\xi))\sum_{j=0}^{K-1}
 \nabla_\theta\log\pi_\theta(\zeta_j\mid x_j,c,\kappa_j)
 \,\middle|\,\xi\right],
 \label{eq:tr-score}
\end{equation}
where $A(B)=\|B-Q\|_{\mathcal W}^2$ and the baseline $b_{\mathrm{base}}$ does not
depend on the sampled actions. This follows by differentiating the
trajectory density; the expected conditional score is zero. It does not
differentiate the physical transition map. A clipped PPO surrogate over
multiple epochs is not the exact finite-change identity in
\eqref{eq:tr-score}, nor a guarantee of improvement in the unshaped task.
In particular, \eqref{eq:tr-variance} retains a variance penalty even when
$Q=\mu_\theta$.

\begin{proposition}[Local task interpretation of response error]
\label{prop:tr-taskerror}
Let $V$ be twice continuously differentiable on a convex neighborhood
$\mathcal N$ containing $q,b$, with
$\|\nabla^2V(y)\|_2\le H$ for all $y\in\mathcal N$. Define
$v=\nabla V(q)$ and $W$ by \eqref{eq:tr-metric}. Then
\begin{equation}
 |V(b)-V(q)|
 \le\sqrt{v^TW^{-1}v}\,\|b-q\|_W
       +\frac H2\|b-q\|_2^2.
 \label{eq:tr-taskerror}
\end{equation}
If $B\in\mathcal N$ almost surely and the required moments are finite,
\begin{equation}
 \mathbb E|V(B)-V(q)|
 \le\sqrt{v^TW^{-1}v}\sqrt{\mathbb E\|B-q\|_W^2}
       +\frac H2\mathbb E\|B-q\|_2^2.
 \label{eq:tr-taskerror-mean}
\end{equation}
\end{proposition}
\begin{proof}
Taylor's formula gives $V(b)-V(q)=v^T(b-q)+r$, with
$|r|\le H\|b-q\|_2^2/2$. Weighted Cauchy--Schwarz gives
$|v^Te|\le\sqrt{v^TW^{-1}v}\sqrt{e^TWe}$. This proves
\eqref{eq:tr-taskerror}. Taking expectations and using
$\mathbb E\sqrt X\le\sqrt{\mathbb EX}$ for $X\ge0$ proves
\eqref{eq:tr-taskerror-mean}.
\end{proof}

If the metric instead uses an approximate gradient $\tilde v$, the
same argument adds $\|v-\tilde v\|_2\|b-q\|_2$ to the bound with
$\tilde v$ in its first term. Further residuals are required when the
observable response does not determine the true continuation task, or
when a frozen-command model continuation is substituted for actual
receding-horizon execution. The omitted latent-memory direction, encoder
reset discrepancies, branch changes, and discrete completion events are
not controlled by this proposition.
Specifically, if $C_{\mathrm{real}}=V(b)+r_{\mathrm{real}}$ and
$C_g=V(q)+r_g$, a bound
$|r_{\mathrm{real}}-r_g|\le\varepsilon_{\mathrm{repr}}$ must be added
separately. Writing this decomposition does not establish that
$\varepsilon_{\mathrm{repr}}$ is small.

\begin{proposition}[Finite-candidate planning comparison]
\label{prop:tr-planning}
Fix a physical policy, a common initial/task/disturbance distribution,
and a nonempty finite command-candidate set $\mathcal C$. Suppose the
actual expected candidate cost $Q$ and model cost $\hat Q$ satisfy
$\sup_{C\in\mathcal C}|Q(C)-\hat Q(C)|\le\varepsilon$. If
$\hat Q(\hat C)\le\min_{C\in\mathcal C}\hat Q(C)+\delta$, then
\begin{equation}
 Q(\hat C)-\min_{C\in\mathcal C}Q(C)\le2\varepsilon+\delta.
 \label{eq:tr-planning}
\end{equation}
\end{proposition}
\begin{proof}
For $C^*\in\arg\min_{C\in\mathcal C}Q(C)$,
\begin{align*}
 Q(\hat C)&\le\hat Q(\hat C)+\varepsilon\\
 &\le\hat Q(C^*)+\delta+\varepsilon\\
 &\le Q(C^*)+\delta+2\varepsilon.
\end{align*}
\end{proof}
The result concerns a common finite candidate set and a fixed physical
policy. It establishes neither continuous-space global CEM optimality
nor actual state-constraint feasibility. It is not a full-task bound for
an evolving policy and a sequence of replanning decisions.

\begin{proposition}[Conditional full-task pair improvement]
\label{prop:tr-pair}
For each $i\in\{-,+\}$, let $\mathcal K_i$ be the complete controller
consisting of its policy and MPC with its paired response model. Use a
common evaluation distribution, horizon, completion rule, and task-cost
definition. If actual and corresponding complete model-evaluation costs
satisfy $|J_i^{\mathrm{real}}-\hat J_i|\le B_i$, then
\begin{equation}
 J_+^{\mathrm{real}}-J_-^{\mathrm{real}}
 \le \hat J_+-\hat J_-+B_++B_-.
 \label{eq:tr-pair}
\end{equation}
Thus $\hat J_--\hat J_+>B_-+B_+$ implies actual task improvement.
\end{proposition}
\begin{proof}
Use $J_+^{\mathrm{real}}\le\hat J_++B_+$ and
$J_-^{\mathrm{real}}\ge\hat J_--B_-$ and subtract.
\end{proof}
The model evaluation in Proposition~\ref{prop:tr-pair} must describe the
same complete replanning controller. Using the training surrogate for a
single frozen remaining command plan requires an additional discrepancy
term. No numerical $B_i$ has been established here. Stochastic-policy
training and deterministic-mean deployment require separate controller
definitions, calibration, and bounds; nonlinear dynamics do not commute
with policy averaging.

\paragraph{Scope and prior work.}
These arguments are standard identities and conditional comparisons.
SPlaTES already couples predictable skills and planning; ReCo already
combines response shaping and policy-aware MPC; guided policy search
already imposes bidirectional policy--trajectory consistency;
costate-focused RL and TaskMet already use task-sensitive model losses.
The candidate must therefore demonstrate its value in task performance,
response predictability, or finite data/computation, rather than claim a
new gradient or optimizer. An ordinary optimizer with the same soft
objective, information, variables, data, derivatives, and constraints may
produce the same update; under the same parameterization and update
rules it can be identical. Command-only trajectory optimization does not
stand in for a comparator allowed to update policy parameters. The soft
penalty alone establishes no independent method advantage.
